\section{Experiential Learning in Engineering Education}\label{experiential-learning-in-engineering-education}

\stanceinfo{\textbf{CFES 2024 Fall General Assembly} [2024-10-27]}{2024-10-27}

\officialstance{The Canadian Federation of Engineering Students (CFES) believes that engineering education should be hands-on and engaging to prepare future engineers for their upcoming careers. The Canadian Federation of Engineering Students further believes that it is essential to expose students to experiential learning through core courses throughout the degree and that institutions should be taking student input on experiential learning.}

\stanceheading{The Students' Position}
\begin{itemize}
 \item The level of experiential learning currently available to Canadian engineering students is heavily dependent on the school being attended, and the province in which the school is located.
 \item Although students can gain experiential learning through internships and co-ops, these opportunities are not guaranteed, nor all-encompassing. It would thus be beneficial to increase the availability of experiential learning offered throughout academic semesters by making coursework more practical and introducing designated environments in which students can develop hands-on skills.
 \item In 2023, the Canadian Engineering Accreditation Board (CEAB) updated their \href{https://engineerscanada.ca/sites/default/files/2023-12/Accreditation_Criteria_Procedures_2023.pdf}{Accreditation Criteria \& Procedures} to add new criteria of 3.4.4.4 as listed below:
 \begin{itemize}
  \item A minimum of 225 AU of engineering design curriculum content in an engineering program shall be delivered by faculty members holding professional engineering licensure as specified in the Interpretive statement on licensure expectations and requirements [CITE].
  \item The CFES agrees with the new criteria added by the CEAB as these faculty members who hold professional engineering licensure will improve the availability and quality of experiential learning opportunities to students that will prepare students to enter the workforce upon graduation.
 \end{itemize}
\end{itemize}

\stanceheading{Recommendations to Partners, Stakeholders, and Other Entities}
\begin{itemize}
 \item The CFES calls on the CEAB to update its accreditation visit practices to incorporate and gather more feedback from students regarding experiential learning, including hands-on aspects of core courses.
 \item The CFES calls on CEAB to create more requirement specifications on the amount of experiential learning required for an engineering degree.
 \item The CFES calls on the Engineering Deans Canada (EDC) to advocate to its member faculties to better construct future curriculums to have a balance of theory and practical skills within the core courses that engineering students are required to take.
 \item The CFES calls on the EDC to direct its member faculties to provide various supports, such as student clubs, to promote experiential learning opportunities along with the faculty.
 \item The CFES further calls on the EDC to direct its member faculties to compile and summarize the wide range of experiential learning opportunities available to their students through a unified resource page, such as an Experiential Learning Hub. This initiative will further ensure that all engineering students can easily discover and access experiential learning opportunities to enrich their degrees.
 \item The CFES urges the EDC to direct its member faculties to take feedback from students regarding experiential learning opportunities.
 \item The CFES calls on its partnered regional organizations (WESST, ESSCO, QCESO, ACES) to collaborate on advocacy efforts to ensure a balance of theory and practical skills are emphasized in their respective regions.
 \item The CFES calls on engineering student societies to support national advocacy efforts and leverage the CEAB policy update to advocate for the implementation of experiential learning in their program.
\end{itemize}

\stanceheading{What the CFES is doing}
\begin{itemize}
 \item The CFES currently offers the opportunity to take part in the Board of European Students of Technology (BEST) courses through partnership, which allows students to cover a technical topic that could be related to career skills.
 \item The CFES holds an annual Canadian Engineering Competition (CEC), that incorporates innovation, creativity, theory, and practical skills throughout its eight (8) competition streams: Re-engineering, Parliamentary Debate, Junior Design, Senior Design, Programming, Innovative Design, Engineering Consulting, and Engineering Communication.
 \item The CFES organizes case competitions at select conferences, providing students with a platform to refine and showcase their problem-solving skills and apply academic knowledge.
\end{itemize}

\stanceheading{What the CFES plans to do}
\begin{itemize}
 \item The CFES will collect data on experiential learning through the national survey.
 \item The CFES will follow up with relevant stakeholders including EDC, to address curriculum changes to increase opportunities for experiential learning in Canadian engineering faculties.
 \item The CFES will emphasize the significance and need for experiential learning during the CEAB and Engineers Canada board meetings.
 \item The CFES will maintain ongoing communications with the International Federation of Engineering Education Societies (IFEES) regarding ways to better advocate for experiential learning opportunities and exchange programs offered to Canadian engineering students.
 \item The CFES will collaborate with the Canadian Engineering Education Association (CEEA) to improve the quality and quantity of experiential learning opportunities in member schools.
\end{itemize}

\newpage
